\documentclass[12pt,a4paper]{article}

\usepackage[T1]{fontenc}
\usepackage{geometry}
\usepackage{hyperref}
\usepackage{booktabs}
\usepackage{tabularx}
\usepackage{longtable}
\usepackage{array}
\usepackage{enumitem}
\usepackage{xcolor}
\usepackage{amsmath}
\usepackage{amssymb}
\usepackage{listings}
\usepackage{parskip}
\usepackage{microtype}
\usepackage{multirow}
\usepackage{mdframed}

\geometry{margin=2.5cm}

\hypersetup{
    colorlinks=true,
    linkcolor=blue!70!black,
    citecolor=green!50!black,
    urlcolor=blue!70!black,
}

\lstset{
    basicstyle=\ttfamily\small,
    frame=single,
    breaklines=true,
    columns=fullflexible,
}

% Coloured boxes for the four analysis dimensions
\newmdenv[backgroundcolor=blue!5,linecolor=blue!40,linewidth=0.8pt,
          innerleftmargin=6pt,innerrightmargin=6pt,innertopmargin=4pt,
          innerbottommargin=4pt]{effbox}
\newmdenv[backgroundcolor=green!6,linecolor=green!40,linewidth=0.8pt,
          innerleftmargin=6pt,innerrightmargin=6pt,innertopmargin=4pt,
          innerbottommargin=4pt]{implbox}
\newmdenv[backgroundcolor=teal!5,linecolor=teal!40,linewidth=0.8pt,
          innerleftmargin=6pt,innerrightmargin=6pt,innertopmargin=4pt,
          innerbottommargin=4pt]{prosbox}
\newmdenv[backgroundcolor=red!5,linecolor=red!30,linewidth=0.8pt,
          innerleftmargin=6pt,innerrightmargin=6pt,innertopmargin=4pt,
          innerbottommargin=4pt]{consbox}

\title{\textbf{Inference Methods for KG-Commit}\\
\large A Survey of Graph Neural Networks, Knowledge Graph Reasoning,\\
and Downstream Prediction Approaches}
\author{KG-Commit Research Project}
\date{\today}

\begin{document}

\maketitle

\noindent\textbf{Legend for shaded analysis boxes:}
\begin{itemize}[noitemsep,leftmargin=1.5em]
  \item \colorbox{blue!5}{\textbf{Efficiency}} --- computational and memory cost for training and inference
  \item \colorbox{green!6}{\textbf{Prior Implementations}} --- readiness and difficulty of adoption
  \item \colorbox{teal!5}{\textbf{Pros}} --- advantages specific to KG-Commit
  \item \colorbox{red!5}{\textbf{Cons}} --- disadvantages and risks specific to KG-Commit
\end{itemize}

\tableofcontents
\newpage

%%=============================================================================
\section{Introduction}
%%=============================================================================

KG-Commit models repository evolution as an incremental heterogeneous knowledge
graph updated with every arriving commit. After the \emph{construction phase}
(building the KG from commits, code subgraphs, authors, issues, and temporal
nodes), the \emph{inference phase} applies reasoning and prediction algorithms
over the populated graph. This document surveys every known inference family
applicable to KG-Commit, from classical graph neural networks to knowledge graph
embedding methods, path-based reasoning, neural-symbolic approaches, temporal
inference, and LLM-augmented methods.

The central inference task is \textbf{Just-In-Time defect prediction}: given the
KG state at commit $c$, predict whether $c$ introduces a defect before the tests
or reviews reveal it. Secondary tasks include \textbf{defect localisation}
(which specific file, class, or function?), \textbf{link prediction} (will this
commit fix issue \#$k$?), \textbf{anomaly detection} (does this commit break an
established pattern?), and \textbf{temporal prediction} (when will the next bug
cluster appear in component $X$?).

For each inference family, four analysis dimensions are evaluated:

\begin{enumerate}[noitemsep]
  \item \textbf{Efficiency} --- training time, inference latency, and memory footprint
  \item \textbf{Prior Implementations and Difficulty} --- tool readiness and integration effort
  \item \textbf{Pros} specific to KG-Commit
  \item \textbf{Cons} specific to KG-Commit
\end{enumerate}

\subsection{The Inference Landscape}

\begin{lstlisting}
KG-Commit Graph
  Node types: COMMIT, FILE, CLASS, FUNCTION, AUTHOR, TIME,
              INTERVAL, BRANCH, ISSUE, STATEMENT, TOKEN, ...
  Edge types: MODIFIES, CALLS, FIXES, CO_CHANGES, DDG_REACHES,
              CFG_TRUE, AUTHORED_BY, PART_OF, ...

Inference Families:
  GNN-based        -- aggregates structural neighbourhoods
  KGE-based        -- embeds entities and relations as vectors/rotations
  Path/Rule-based  -- reasons over multi-hop paths and logical rules
  Temporal         -- exploits commit sequence and timestamps
  Subgraph         -- inductive reasoning from local subgraph patterns
  LLM-augmented    -- natural language grounding over KG facts
  Self-supervised  -- learns from unlabelled graph structure
  Graph Transformer-- full self-attention with structural encodings
\end{lstlisting}

\subsection{Key Principle}

The KG-Commit graph is \emph{heterogeneous} (multiple node/edge types),
\emph{temporal} (commits arrive in sequence), \emph{incremental} (updated
per commit), and \emph{label-scarce} (most commits are unlabelled or clean).
Any inference method must handle all four properties well to be practically
deployable.

\subsection{General Design Axes}

\begin{longtable}{p{3.2cm}p{5cm}p{4.5cm}}
\toprule
\textbf{Axis} & \textbf{Question} & \textbf{Why it matters} \\
\midrule
\endhead
Transductive vs.\ inductive & Fixed graph vs.\ new nodes at test time & New repos and projects require inductive methods \\
Homogeneous vs.\ heterogeneous & Single type vs.\ multiple node/edge types & KG-Commit has $>$15 node types and $>$30 edge types \\
Shallow vs.\ deep & 1--2 hops vs.\ multi-hop reasoning & Defects often span multiple hops \\
Symbolic vs.\ neural & Interpretable rules vs.\ dense vectors & Both wanted: explainability and accuracy \\
Temporal awareness & Ignores timestamps vs.\ models commit sequence & Bug patterns have strong temporal locality \\
Scalability & Linear in graph size vs.\ quadratic & Full Apache history = $10^7$+ nodes \\
Label efficiency & Supervised vs.\ self-supervised & JIT defect labels are scarce ($<$5\% positive) \\
\bottomrule
\end{longtable}

%%=============================================================================
\section{Family 1: Shallow Graph Convolutional Networks}
%%=============================================================================

\subsection{What It Captures}

Spectral and spatial graph convolutions aggregate information from $k$-hop
neighbourhoods via learnable weight matrices applied to the normalised adjacency.
The canonical model, GCN (Kipf \& Welling, 2017), computes:
\[
  H^{(\ell+1)} = \sigma\!\left(\tilde{D}^{-\frac{1}{2}}\tilde{A}\tilde{D}^{-\frac{1}{2}}H^{(\ell)}W^{(\ell)}\right)
\]
where $\tilde{A} = A + I$ is the adjacency with self-loops. GraphSAGE (Hamilton
et al., 2017) samples a fixed-size neighbourhood and concatenates the aggregated
neighbour vector with the central node, enabling inductive generalisation to
unseen nodes. ChebNet (Defferrard et al., 2016) uses Chebyshev polynomial filters
of order $K$, approximating multi-hop spectral convolution. SGC (Wu et al., 2019)
removes non-linearities between layers, reducing $K$-layer GCN to a single linear
transform of the $K$-th power of the adjacency.

\paragraph{Model formulation for KG-Commit.}
\begin{lstlisting}
Input node features:
  COMMIT   [la, ld, nf, fix, ns, nd, entropy, ...]
  FUNCTION [loc, cyclomatic, fan_in, fan_out, ...]
Aggregation: mean/sum/max over 1-hop neighbours per edge type
Output: per-node embedding -> defect prediction head (MLP + sigmoid)
Baseline: treat all edges as a single type (homogeneous projection)
\end{lstlisting}

\begin{effbox}
\textbf{Efficiency.}
Training: $O(|\mathcal{E}| \cdot d^2)$ per layer for $d$-dimensional features
and $|\mathcal{E}|$ edges. A 2-layer GCN over the KG-Commit graph with $10^5$
nodes and $10^6$ edges trains in minutes on a single GPU. Inference: one forward
pass per commit subgraph, $O(k \cdot d^2)$ for $k$ neighbours.
Mini-batch training via neighbourhood sampling (GraphSAGE) reduces memory to
$O(\text{batch} \times k^L \cdot d)$ for $L$ layers. SGC inference is a single
sparse matrix multiply --- extremely fast. The most compute-efficient GNN family.
\end{effbox}

\begin{implbox}
\textbf{Prior Implementations and Difficulty.}
\textbf{PyTorch Geometric} and \textbf{DGL} provide production-ready GCN,
GraphSAGE, ChebNet, and SGC implementations. Both support heterogeneous graphs
via wrapper modules. \textbf{Stellargraph} (TensorFlow-based) provides additional
baselines. Integration with KG-Commit: export to a PyG \texttt{HeteroData} object
(via py2neo or neo4j-arrow). No graph-type-specific code required for a
homogeneous baseline. Difficulty: \textbf{very easy}.
\end{implbox}

\begin{prosbox}
\textbf{Pros.}
Simplest and fastest GNN baseline; essential for ablation studies. GraphSAGE
inductive mode handles new repositories and projects at test time without
retraining. SGC's linear formulation is interpretable --- the weight matrix
directly shows which node features drive predictions. JITLine (2022) uses GCN
over file co-change graphs as a competitive JIT-DP baseline, validating this
family for the task.
\end{prosbox}

\begin{consbox}
\textbf{Cons.}
Homogeneous GCN ignores the rich edge-type vocabulary of KG-Commit: all edges
are treated equally, meaning \texttt{MODIFIES} and \texttt{CO\_CHANGES\_WITH}
carry equal weight. GraphSAGE with uniform sampling may under-sample rare but
critical edge types (e.g., \texttt{FIXES} edges to issue nodes). Over-smoothing
occurs at $\geq$4 layers, washing out local defect signals. No temporal awareness.
\end{consbox}

\paragraph{Key papers.} GCN (Kipf \& Welling, ICLR 2017), GraphSAGE (Hamilton
et al., NeurIPS 2017), ChebNet (Defferrard et al., NeurIPS 2016), SGC (Wu et al.,
ICML 2019).

%%=============================================================================
\section{Family 2: Attention-Based Graph Neural Networks}
%%=============================================================================

\subsection{What It Captures}

Graph Attention Networks (GAT, Veli\v{c}kovi\'{c} et al., 2018) learn
importance weights over neighbours via a learnable attention mechanism:
\[
  \alpha_{ij} = \frac{\exp(\text{LeakyReLU}(\mathbf{a}^\top [\mathbf{W}h_i \| \mathbf{W}h_j]))}
                     {\sum_{k \in \mathcal{N}(i)} \exp(\text{LeakyReLU}(\mathbf{a}^\top [\mathbf{W}h_i \| \mathbf{W}h_k]))}
\]
GATv2 (Brody et al., 2022) fixes expressiveness limitations of GAT by applying
the linear transform \emph{after} concatenation. The Heterogeneous Attention
Network (HAN, Wang et al., 2019) extends this to heterogeneous graphs via
meta-path-based attention: first aggregate within each meta-path, then attend
over meta-paths with a semantic-level attention vector.

\paragraph{Model formulation for KG-Commit.}
\begin{lstlisting}
Meta-paths for HAN:
  COMMIT --MODIFIES--> FILE --MODIFIES<-- COMMIT    (co-edit similarity)
  COMMIT --AUTHORED_BY--> AUTHOR --AUTHORED_BY<-- COMMIT  (co-author)
  FUNCTION --CALLS--> FUNCTION --CALLS<-- FUNCTION  (call-site coupling)
  COMMIT --FIXES--> ISSUE <--MENTIONS-- FUNCTION    (issue-function link)
Node-level attention: weight neighbours by edge type and feature similarity
\end{lstlisting}

\begin{effbox}
\textbf{Efficiency.}
GAT: $O(|\mathcal{E}| \cdot d^2 + |\mathcal{E}| \cdot K)$ per layer for $K$
attention heads. In practice $2$--$3\times$ slower than GCN per epoch.
GATv2: same asymptotic cost, slightly higher constant.
HAN: meta-path materialisation adds a preprocessing step ($O(|\mathcal{E}|)$
per short meta-path). For KG-Commit with $10^6$ edges: HAN preprocessing takes
1--5 minutes; GAT training 10--30 minutes per epoch on GPU.
\end{effbox}

\begin{implbox}
\textbf{Prior Implementations and Difficulty.}
\textbf{PyTorch Geometric}: \texttt{GATConv}, \texttt{GATv2Conv},
\texttt{HANConv}. \textbf{DGL}: \texttt{GATConv}, \texttt{HAN}.
Meta-path definition is a configuration file listing meta-path schemas,
which must be manually specified for KG-Commit's schema (a one-time design
task). Meta-path materialisation is trivially computable via Cypher queries
on the Neo4j graph. GATv2 should be preferred over GAT for expressiveness.
Difficulty: GAT/GATv2 \textbf{very easy}; HAN (with meta-path setup)
\textbf{easy}.
\end{implbox}

\begin{prosbox}
\textbf{Pros.}
Attention weights are \emph{interpretable}: which neighbours most influenced
a defect prediction can be extracted and shown to developers. Multi-head
attention naturally captures multiple roles of a commit (large change AND
authored by a low-ownership author). HAN's meta-path design aligns naturally
with the KG-Commit schema and encodes domain knowledge about which structural
patterns predict defects.
\end{prosbox}

\begin{consbox}
\textbf{Cons.}
HAN is sensitive to the choice of meta-paths; wrong meta-paths hurt performance.
Meta-path materialisation can be expensive for long paths in a large KG.
Standard GAT is provably no more expressive than GCN for certain graph
structures --- motivation for GATv2. Attention weights are not globally
normalised across the whole graph, limiting their use as a global saliency
measure. No relational weight separation (all edge types share the same $W$).
\end{consbox}

\paragraph{Key papers.} GAT (Veli\v{c}kovi\'{c} et al., ICLR 2018), GATv2
(Brody et al., ICLR 2022), HAN (Wang et al., WWW 2019).

%%=============================================================================
\section{Family 3: Gated and Sequential Graph Neural Networks}
%%=============================================================================

\subsection{What It Captures}

The Gated Graph Neural Network (GGNN, Li et al., 2016) replaces the standard
GNN update rule with a Gated Recurrent Unit (GRU) applied over $T$ propagation
steps:
\[
  h_v^{(t)} = \text{GRU}\!\left(h_v^{(t-1)},\;\text{AGG}\{h_u^{(t-1)}: u \in \mathcal{N}(v)\}\right)
\]
This allows information to flow over longer paths without over-smoothing for a
fixed $T$. Tree-LSTM (Tai et al., 2015) generalises the LSTM cell to tree
structures, applicable to the AST subtree of a function. GGNN is the backbone
of \textbf{Devign} (Zhou et al., 2019), the most widely cited GNN-based
vulnerability detection system, applied directly to CPG subgraphs.

\paragraph{Model formulation for KG-Commit.}
\begin{lstlisting}
GGNN over CPG subgraph (Devign-style):
  Initialise: node embedding from token/type features
  Propagate T=8 steps: GRU update over outgoing CPG edges
  Readout: gated sum aggregation over all node embeddings
  Output: defect probability (sigmoid)

Tree-LSTM over AST:
  Leaf-to-root information flow for syntactic features
  Applicable to the AST subgraphs produced by tree-sitter / Spoon
\end{lstlisting}

\begin{effbox}
\textbf{Efficiency.}
GGNN: $O(T \cdot |\mathcal{E}| \cdot d^2)$ for $T$ propagation steps. In Devign,
$T=8$ and $d=200$ is standard. For a single CPG subgraph (500--2000 nodes), a
forward pass takes $<$10ms on GPU. Batched over 100 commits: $<$1 second.
Training over 10,000 commits: 1--2 hours on GPU. Tree-LSTM: $O(n \cdot d^2)$
per tree for $n$ AST nodes. Both are well within single-GPU memory for
function-level graphs.
\end{effbox}

\begin{implbox}
\textbf{Prior Implementations and Difficulty.}
\textbf{Devign} reference code (available on GitHub) uses a custom PyTorch GGNN.
\textbf{DGL}: \texttt{GatedGraphConv}. \textbf{PyTorch Geometric}:
\texttt{GatedGraphConv}. \textbf{Tree-LSTM}: DGL \texttt{TreeLSTMCell}.
Devign's Java port requires generating CPG data via Joern and adapting the graph
serialisation. The upstream Devign code is research-grade but straightforward to
adapt. Difficulty: GGNN \textbf{easy}; full Devign pipeline on Java CPG
\textbf{moderate}.
\end{implbox}

\begin{prosbox}
\textbf{Pros.}
GGNN is the academically canonical choice for code-graph defect detection
(Devign, ReVeal, dozens of follow-ups). GRU gating prevents over-smoothing
even for $T=8$ steps, enabling long-range information propagation over the CPG.
Tree-LSTM is uniquely suited to the hierarchical AST structure built in the
construction phase. Both are directly applicable to the CPG subgraphs produced
by Joern in the KG-Commit pipeline.
\end{prosbox}

\begin{consbox}
\textbf{Cons.}
Fixed propagation step count $T$ must be tuned; for very deep CPGs $T$ may
need to be large (slow). Tree-LSTM is only applicable to the AST subgraph
(no cycles), not the full CPG. Sequential GRU recurrence prevents full GPU
parallelisation across time steps, making GGNN slower per epoch than GCN.
Does not natively handle multiple edge types without type-specific adjacency
matrices.
\end{consbox}

\paragraph{Key papers.} GGNN (Li et al., ICLR 2016), Devign (Zhou et al.,
NeurIPS 2019), Tree-LSTM (Tai et al., ACL 2015), ReVeal (Chakraborty et al., 2022).

%%=============================================================================
\section{Family 4: Relational and Heterogeneous Graph Neural Networks}
\label{sec:rgcn}
%%=============================================================================

\subsection{What It Captures}

Relational GCN (R-GCN, Schlichtkrull et al., 2018) uses separate weight matrices
per relation (edge type):
\[
  h_i^{(\ell+1)} = \sigma\!\left(\sum_{r \in \mathcal{R}} \sum_{j \in \mathcal{N}_r(i)} \frac{1}{c_{i,r}} W_r^{(\ell)} h_j^{(\ell)} + W_0^{(\ell)} h_i^{(\ell)}\right)
\]
This directly models the distinct semantics of each edge type in KG-Commit.
The Heterogeneous Graph Transformer (HGT, Hu et al., 2020) extends this with a
transformer-style attention mechanism across node and edge types, using
type-specific linear projections for keys, queries, and values. CompGCN
(Vashishth et al., 2020) treats entities and relations as co-embedded objects,
enabling composition operations (subtract, multiply, circular correlation) over
$(h, r, t)$ triples --- bridging GNN and KGE families.

\paragraph{Model formulation for KG-Commit.}
\begin{lstlisting}
R-GCN / HGT over KG-Commit:
  Node types : COMMIT, FILE, FUNCTION, AUTHOR, ISSUE, CLASS, ...
  Edge types  : MODIFIES, CALLS, FIXES, DDG_REACHES, CFG_TRUE,
                AUTHORED_BY, CO_CHANGES, HAS_METRIC, ...
  Per-relation weight matrix W_r (R-GCN, with basis decomposition)
  Per-(src_type, rel, dst_type) transformer head (HGT)
  Prediction head: sigmoid over COMMIT node embedding
\end{lstlisting}

\begin{effbox}
\textbf{Efficiency.}
R-GCN: $O(|\mathcal{R}| \cdot |\mathcal{E}| \cdot d^2)$. With basis
decomposition (R-GCN option), reduced to $O(B \cdot d^2 + |\mathcal{R}|
\cdot B \cdot d)$ for $B$ basis matrices ($B \ll |\mathcal{R}|$).
For KG-Commit with $\approx$30 edge types, basis decomposition is essential.
HGT: $O(|\mathcal{E}| \cdot d^2)$ with a higher constant than GCN due to
transformer heads. Both are 2--5$\times$ slower than GCN per epoch but
dramatically more accurate on heterogeneous KGs.
\end{effbox}

\begin{implbox}
\textbf{Prior Implementations and Difficulty.}
\textbf{PyTorch Geometric}: \texttt{RGCNConv} (with basis decomposition),
\texttt{HGTConv}, \texttt{RGATConv}, \texttt{CompGCN}.
\textbf{DGL}: \texttt{RelGraphConv}, \texttt{HGTConv}.
Both libraries natively support \texttt{HeteroData} objects with typed
node/edge features. Integration: export KG-Commit to a PyG \texttt{HeteroData}
object from Neo4j. R-GCN basis decomposition should be used when
$|\mathcal{R}| > 15$. Difficulty: R-GCN \textbf{easy}; HGT \textbf{easy};
CompGCN \textbf{moderate} (requires relation embedding design).
\end{implbox}

\begin{prosbox}
\textbf{Pros.}
The \textbf{single most appropriate GNN family} for KG-Commit. The
heterogeneous edge vocabulary (MODIFIES vs.\ CALLS vs.\ FIXES) encodes
fundamentally different semantics that R-GCN/HGT treats correctly.
HGT's transformer attention over node types naturally handles the size
disparity between FUNCTION nodes (many, small features) and AUTHOR nodes
(few, rich historical features). CompGCN's composition operations encode
relational reasoning (``commit by low-ownership author who modified
high-churn file'') as a composed entity embedding.
\end{prosbox}

\begin{consbox}
\textbf{Cons.}
R-GCN parameter count grows with $|\mathcal{R}|$ (mitigated by basis
decomposition). HGT requires type-specific feature projections, adding
preprocessing complexity when node feature dimensionalities differ across
types. Both require the full heterogeneous graph to be in memory or carefully
mini-batched. HGT is harder to interpret than GAT (attention weights span
types and are not directly human-readable).
\end{consbox}

\paragraph{Key papers.} R-GCN (Schlichtkrull et al., ESWC 2018), HGT
(Hu et al., WWW 2020), CompGCN (Vashishth et al., ICLR 2020).

%%=============================================================================
\section{Family 5: Translational Knowledge Graph Embeddings}
%%=============================================================================

\subsection{What It Captures}

Translational KGE methods embed entities $h, t$ and relations $r$ in a
low-dimensional vector space such that $\mathbf{h} + \mathbf{r} \approx \mathbf{t}$
for observed triples $(h, r, t)$. TransE (Bordes et al., 2013) is the canonical
model. Variants address its limitations: TransH (Wang et al., 2014) projects
entities onto relation-specific hyperplanes, enabling many-to-many relations.
TransR (Lin et al., 2015) uses relation-specific projection matrices.
TransD (Ji et al., 2015) uses entity-relation-specific dynamic mappings.

\paragraph{Scoring functions.}
\begin{lstlisting}
TransE:  f(h,r,t) = -||h + r - t||
TransH:  f(h,r,t) = -||h_perp + d_r - t_perp||
TransR:  f(h,r,t) = -||M_r h + r - M_r t||
Training: margin-based ranking loss with negative sampling
Inference: score all candidate tails t' given (h, r, ?)
KG-Commit triples: (COMMIT, MODIFIES, FILE),
                   (COMMIT, AUTHORED_BY, AUTHOR), ...
\end{lstlisting}

\begin{effbox}
\textbf{Efficiency.}
Training: $O(|\mathcal{T}| \cdot d)$ per epoch for $|\mathcal{T}|$ triples
and $d$-dimensional embeddings (typically $d=100$--$300$). Extremely fast:
$10^6$ triples train in minutes on CPU. Inference: $O(N \cdot d)$ per query
(score all $N$ entities as tail candidates). Total model size: $O((N+R) \cdot d)$
--- typically a few MB. The most memory-efficient inference family. Ideal for
CPU-only deployment or resource-constrained environments.
\end{effbox}

\begin{implbox}
\textbf{Prior Implementations and Difficulty.}
\textbf{PyKEEN} (Fraunhofer, production-ready) implements TransE, TransH,
TransR, TransD, and dozens of others in a unified framework with automated
hyperparameter optimisation. \textbf{OpenKE} (THU, production-ready) provides
C++-accelerated training. KG-Commit triple extraction: export a triple file
from Neo4j via Cypher in one query. Difficulty: \textbf{very easy}.
\end{implbox}

\begin{prosbox}
\textbf{Pros.}
Zero GNN infrastructure required. CPU-deployable for online inference after
training. The triple-scoring formulation maps directly to link prediction
queries: ``predict which ISSUE this COMMIT will FIX,'' ``predict which
FUNCTION will be changed next.'' Highly scalable: training on the full
KG-Commit triple set ($10^7$ triples) is feasible on a single machine.
\end{prosbox}

\begin{consbox}
\textbf{Cons.}
TransE cannot model symmetric or one-to-many relations well (e.g., one AUTHOR
modifies many FILEs). Purely structural --- no code features incorporated.
The translation assumption does not capture multi-hop compositionality without
explicit path encoding. All KGE models are transductive: embedding new commits
at test time requires re-training or a separate inductive wrapper. Lower accuracy
than GNNs on feature-rich graphs.
\end{consbox}

\paragraph{Key papers.} TransE (Bordes et al., NeurIPS 2013), TransH (Wang et al.,
AAAI 2014), TransR (Lin et al., AAAI 2015), PyKEEN (Ali et al., JMLR 2021).

%%=============================================================================
\section{Family 6: Bilinear Knowledge Graph Embeddings}
%%=============================================================================

\subsection{What It Captures}

Bilinear KGE models score a triple $(h, r, t)$ via a bilinear interaction.
DistMult (Yang et al., 2015) uses a diagonal bilinear form $f = \langle h, r,
t \rangle$. ComplEx (Trouillon et al., 2016) lifts to complex space enabling
asymmetric relations: $f = \text{Re}(\langle h, r, \bar{t} \rangle)$.
TuckER (Bala\v{z}evi\'{c} et al., 2019) uses Tucker decomposition over a core
tensor $\mathcal{W}$: $f = \mathcal{W} \times_1 h \times_2 r \times_3 t$.
RESCAL (Nickel et al., 2011) uses a full rank-$d$ bilinear form $f = h^\top
M_r t$, the most expressive but most expensive variant.

\paragraph{Scoring functions.}
\begin{lstlisting}
DistMult:  f(h,r,t) = sum_i h_i * r_i * t_i
ComplEx:   f(h,r,t) = Re(<h, r, conj(t)>)    (complex vectors)
TuckER:    f(h,r,t) = W x_1 h x_2 r x_3 t   (Tucker decomp)
Training:  binary cross-entropy with negative sampling (DistMult, ComplEx)
           or multinomial log-loss (TuckER)
\end{lstlisting}

\begin{effbox}
\textbf{Efficiency.}
DistMult: $O(|\mathcal{T}| \cdot d)$ per epoch --- the cheapest bilinear
model. ComplEx: same asymptotic cost, $2\times$ parameters (real + imaginary).
TuckER: $O(|\mathcal{T}| \cdot d_e \cdot d_r)$ where $d_e, d_r$ are entity
and relation dimensions --- slightly more expensive but achieves SOTA on
standard benchmarks. RESCAL: $O(|\mathcal{T}| \cdot d^2)$ --- quadratic in
$d$, impractical for $d > 100$. For KG-Commit: DistMult or ComplEx for fast
training; TuckER for higher accuracy.
\end{effbox}

\begin{implbox}
\textbf{Prior Implementations and Difficulty.}
\textbf{PyKEEN}: all models above with one-line instantiation. \textbf{AmpliGraph}
(TensorFlow, production-ready) provides DistMult, ComplEx with early stopping and
evaluation metrics. \textbf{LibKGE} (Broscheit et al.) is a research framework
with extensive hyperparameter search. All accept standard triple files.
Difficulty: DistMult/ComplEx \textbf{very easy}; TuckER \textbf{easy};
RESCAL \textbf{easy} (but slow for large $d$).
\end{implbox}

\begin{prosbox}
\textbf{Pros.}
ComplEx handles asymmetric relations correctly: AUTHORED\_BY is asymmetric
(AUTHOR$\to$COMMIT $\neq$ COMMIT$\to$AUTHOR), which TransE cannot model.
DistMult achieves competitive results with minimal compute. TuckER achieves
SOTA link prediction on Freebase and Wikidata benchmarks. The bilinear form
enables matrix-factorisation-style analysis: entity embeddings can be
interpreted as latent ``roles'' in the code change process.
\end{prosbox}

\begin{consbox}
\textbf{Cons.}
DistMult cannot model anti-symmetric relations. No structural awareness of
graph topology beyond the training triples. Cannot incorporate node feature
vectors (code metrics, embeddings) without additional wrapper layers. Bilinear
models tend to underperform GNNs when rich node features are available.
Transductive limitation applies to all variants.
\end{consbox}

\paragraph{Key papers.} RESCAL (Nickel et al., ICML 2011), DistMult (Yang et al.,
ICLR 2015), ComplEx (Trouillon et al., ICML 2016), TuckER (Bala\v{z}evi\'{c}
et al., EMNLP 2019).

%%=============================================================================
\section{Family 7: Rotation-Based Knowledge Graph Embeddings}
%%=============================================================================

\subsection{What It Captures}

RotatE (Sun et al., 2019) models each relation as a rotation in complex space:
$\mathbf{t} = \mathbf{h} \circ \mathbf{r}$, where $\circ$ is element-wise
complex multiplication and $|r_k| = 1$. This unified framework models three
key relation patterns: symmetry, antisymmetry, and composition ($r_3 = r_1
\circ r_2$). HAKE (Zhang et al., 2020) adds a polar-coordinate modulus
component to model hierarchical relations (e.g., FILE $\supset$ CLASS $\supset$
FUNCTION). PairRE (Chao et al., 2021) uses two separate relation vectors
(head projection and tail projection), improving multi-hop chain modelling.
QuatE (Zhang et al., 2019) generalises to quaternion space (4D rotations).

\paragraph{Scoring functions.}
\begin{lstlisting}
RotatE:  f(h,r,t) = -||h circ r - t||,   |r_k|=1  (complex)
HAKE:    f(h,r,t) = modulus term + phase term       (polar coords)
PairRE:  f(h,r,t) = -||h circ r_H - t circ r_T||
QuatE:   f(h,r,t) = (h otimes r) . t               (quaternion product)
\end{lstlisting}

\begin{effbox}
\textbf{Efficiency.}
RotatE: $O(|\mathcal{T}| \cdot d)$ per epoch, same as ComplEx. Self-adversarial
negative sampling (RotatE's key trick) adds a small constant overhead.
HAKE: same asymptotic cost, slightly higher due to dual-component scoring.
PairRE and QuatE: $2\times$ and $4\times$ parameter count respectively, but
similar per-epoch cost. GPU training on $10^7$ triples: 1--2 hours.
CPU inference: $<$1ms per link prediction query.
\end{effbox}

\begin{implbox}
\textbf{Prior Implementations and Difficulty.}
\textbf{PyKEEN}: RotatE, HAKE, PairRE, QuatE all available.
\textbf{RotatE} reference code (Sun et al.) is available on GitHub and is
actively maintained. HAKE reference code is also available. Self-adversarial
negative sampling is the key hyperparameter (temperature $\alpha$).
Difficulty: RotatE/PairRE via PyKEEN \textbf{very easy}; HAKE and QuatE
\textbf{easy}.
\end{implbox}

\begin{prosbox}
\textbf{Pros.}
RotatE is the \textbf{current SOTA single-model KGE} on standard benchmarks
(FB15k-237, WN18RR). HAKE's hierarchical modelling directly fits KG-Commit's
containment hierarchy (REPO $\supset$ FILE $\supset$ CLASS $\supset$ FUNCTION).
PairRE's multi-hop composition models relation chains (COMMIT $\to$ FUNCTION
$\to$ CALLS $\to$ BUG) that appear naturally in defect propagation. Rotation
constraints are geometrically interpretable.
\end{prosbox}

\begin{consbox}
\textbf{Cons.}
Transductive: same limitation as all KGE methods. The rotation constraint
$|r_k|=1$ restricts relation expressiveness for non-rotational semantics.
Self-adversarial negative sampling has a temperature hyperparameter that
requires tuning. QuatE requires $4\times$ embedding dimensions, expensive
for large entity sets. No code feature incorporation without a wrapper.
\end{consbox}

\paragraph{Key papers.} RotatE (Sun et al., ICLR 2019), HAKE (Zhang et al.,
ACL 2020), PairRE (Chao et al., ACL 2021), QuatE (Zhang et al., NeurIPS 2019).

%%=============================================================================
\section{Family 8: Path-Based Reasoning}
%%=============================================================================

\subsection{What It Captures}

Path-based methods reason over multi-hop paths in the KG. The Path Ranking
Algorithm (PRA, Lao et al., 2011) enumerates relation paths between entity
pairs and uses them as features for a logistic regression classifier.
DeepPath (Xiong et al., 2017) trains a reinforcement learning agent to walk
the KG, learning which paths are most predictive. MINERVA (Das et al., 2018)
is a fully end-to-end RL path reasoner that navigates from the query entity
to the answer entity. M-Walk (Shen et al., 2018) combines MCTS with LSTM-guided
path walking for more stable training.

\paragraph{Path formulation for KG-Commit.}
\begin{lstlisting}
Predictive paths for defect prediction:
  COMMIT --MODIFIES--> FILE --HAS_SMELL--> GodClass
  COMMIT --AUTHORED_BY--> AUTHOR --LOW_OWNERSHIP--> FILE
  FUNCTION --CALLS--> FUNCTION --MAY_BE_NULL--> Statement
  COMMIT --CHANGES--> FUNCTION <--BACKWARD_SLICE-- Criterion

PRA: enumerate paths up to length L=3; use as binary features
DeepPath/MINERVA: RL agent navigates from COMMIT to predicted label
\end{lstlisting}

\begin{effbox}
\textbf{Efficiency.}
PRA: path enumeration is $O(|\mathcal{E}|^L)$ for path length $L$, tractable
for $L \leq 3$ with random walk approximation. Feature extraction and logistic
regression: fast. DeepPath/MINERVA: RL training is slow ($10^3$--$10^4$
episodes per relation), typically hours on GPU. Inference: a single forward
pass of the policy network per query $O(L \cdot d)$ --- very fast once trained.
\end{effbox}

\begin{implbox}
\textbf{Prior Implementations and Difficulty.}
PRA: \textbf{ProPPR} (Cohen et al.) provides a production implementation;
random walk approximation is trivially implementable in Python with NetworkX.
\textbf{DeepPath}: reference PyTorch code on GitHub (Xiong et al.).
\textbf{MINERVA}: reference code on GitHub (Das et al.), maintained.
None are currently integrated into any code analysis pipeline. Integration
requires a custom KG-Commit environment wrapper for the RL agent.
Difficulty: PRA \textbf{moderate}; DeepPath/MINERVA \textbf{hard}
(RL training instability is a known challenge with sparse rewards).
\end{implbox}

\begin{prosbox}
\textbf{Pros.}
Highly \textbf{interpretable}: paths like ``COMMIT $\to$ MODIFIES $\to$ FILE
$\to$ HAS\_SMELL $\to$ GodClass'' directly explain why a commit is flagged.
PRA paths can be displayed to developers as human-readable explanations.
Multi-hop reasoning captures change-propagation patterns not visible at
1-hop depth. MINERVA navigates through the actual KG topology, respecting
graph structure and relation semantics.
\end{prosbox}

\begin{consbox}
\textbf{Cons.}
PRA path enumeration scales poorly to long paths ($L \geq 4$) in large KGs.
RL-based methods suffer from sparse reward signals --- defect labels are rare
($<$5\% of commits), making reward assignment difficult. Path-based methods
miss structural patterns not on any single path. M-Walk's MCTS requires
multiple rollouts per training step, amplifying the sparsity problem.
\end{consbox}

\paragraph{Key papers.} PRA (Lao et al., EMNLP 2011), DeepPath (Xiong et al.,
EMNLP 2017), MINERVA (Das et al., ICLR 2018), M-Walk (Shen et al., NeurIPS 2018).

%%=============================================================================
\section{Family 9: Rule Learning and Neural-Symbolic Methods}
%%=============================================================================

\subsection{What It Captures}

Rule learning methods mine logical Horn rules from the KG: e.g.,
$A(x,y) \wedge B(y,z) \Rightarrow C(x,z)$ with a confidence score.
AMIE (Gal\'{a}rraga et al., 2013) mines such rules by exhaustive search with
statistical pruning. Neural-LP (Yang et al., 2017) learns differentiable rule
weights as a soft version of Datalog. DRUM (Sadeghian et al., 2019) mines rules
via bidirectional recurrent networks. RNNLogic (Qu et al., 2021) combines rule
generation (RNN) with rule-guided reasoning (EM algorithm). These methods
produce \emph{symbolic rules} that can be audited by humans and encoded directly
into review checklists.

\paragraph{Rule examples for KG-Commit.}
\begin{lstlisting}
Mined rule 1 (defect predictor):
  MODIFIES(c, f) AND HAS_SMELL(f, GodClass) AND LOC_DELTA > 50
    => LIKELY_BUGGY(c)

Mined rule 2 (incomplete change):
  MODIFIES(c, f1) AND CO_CHANGES(f1, f2) AND NOT MODIFIES(c, f2)
    => INCOMPLETE_CHANGE(c)

Mined rule 3 (ownership risk):
  AUTHORED_BY(c, a) AND OWNERSHIP(a, f) < 0.1 AND MODIFIES(c, f)
    => HIGH_RISK(c)
\end{lstlisting}

\begin{effbox}
\textbf{Efficiency.}
AMIE: rule mining scales as $O(|\mathcal{T}|^2)$ worst case, but with
statistical pruning runs in minutes on $10^6$ triples. Neural-LP/DRUM:
$O(|\mathcal{T}| \cdot d \cdot L)$ per epoch for path length $L$ ---
similar to KGE. RNNLogic: EM alternation, 10--50 iterations each
$O(|\mathcal{T}| \cdot d)$. After training, rule application is
$O(|\mathcal{T}|)$ --- very fast. Rule sets are tiny (dozens to hundreds
of rules), enabling sub-millisecond inference.
\end{effbox}

\begin{implbox}
\textbf{Prior Implementations and Difficulty.}
\textbf{AMIE 3.0}: production-ready Java tool, actively maintained.
\textbf{Neural-LP}: reference TensorFlow code (Yang et al.).
\textbf{DRUM}: PyTorch reference code.
\textbf{RNNLogic}: PyTorch, maintained.
For KG-Commit: AMIE can be run directly on the exported triple file; mined
rules can be reviewed and seeded into the system as interpretable inference
rules. This is a strong approach for building \emph{auditable defect
prediction rules}. Difficulty: AMIE \textbf{easy}; Neural-LP/DRUM
\textbf{moderate}; RNNLogic \textbf{moderate}.
\end{implbox}

\begin{prosbox}
\textbf{Pros.}
\textbf{Fully interpretable}: mined rules are auditable by domain experts.
AMIE can discover novel co-change rules, smell-defect correlations, and
ownership-risk patterns not encoded a priori in the KG schema. Rules are
directly actionable as code review checklists. Neural-LP rules are
differentiable, enabling end-to-end training. Symbolic rules are portable
across projects without retraining on new code.
\end{prosbox}

\begin{consbox}
\textbf{Cons.}
Rule mining requires a sufficient number of positive examples to achieve
statistical confidence --- scarce with few defect-labelled commits.
Mined rules may be spurious correlations (not causal). Neural-LP and DRUM
do not scale as well as GNNs on feature-rich graphs. Rules mined from one
project may not transfer to another (different coding style, different bug
classes). No node feature incorporation in purely symbolic rule mining.
\end{consbox}

\paragraph{Key papers.} AMIE (Gal\'{a}rraga et al., WWW 2013; AMIE 3, VLDB 2020),
Neural-LP (Yang et al., NeurIPS 2017), DRUM (Sadeghian et al., NeurIPS 2019),
RNNLogic (Qu et al., ICLR 2021).

%%=============================================================================
\section{Family 10: Temporal Knowledge Graph Inference}
%%=============================================================================

\subsection{What It Captures}

Temporal KGE methods extend entity and relation embeddings with a time component.
TTransE (Leblay \& Chekol, 2018) adds a time translation: $\mathbf{h} + \mathbf{r}
+ \boldsymbol{\tau}(t) \approx \mathbf{tail}$. TNTComplEx (Lacroix et al., 2020)
uses a 4th-order tensor decomposition over (entity, relation, entity, time).
RE-NET (Jin et al., 2020) models relational event sequences with a GRU over
temporal subgraphs. TiRGN (Li et al., 2022) adds a recurrent encoder over
historical entity patterns combined with a local graph encoder. TANGO (Han et al.,
2021) uses Neural ODEs to model continuous temporal evolution of entity embeddings.

\paragraph{Temporal formulation for KG-Commit.}
\begin{lstlisting}
Temporal triples:
  (COMMIT_t, MODIFIES, FUNCTION, timestamp_t)

Prediction: given history up to t-1, predict if COMMIT_t is buggy

TTransE:  h + r + tau(t) ~= tail  (tau = time embedding lookup)
RE-NET:   GRU over sequence of historical commit subgraphs
TANGO:    dh/dt = f(h, t)  via Neural ODE (continuous evolution)

KG-Commit structural assets: TIME nodes, INTERVAL nodes,
  timestamp edges, author commit frequency over time windows
\end{lstlisting}

\begin{effbox}
\textbf{Efficiency.}
TTransE/TNTComplEx: same cost as static counterparts plus time embedding
storage $O(T \cdot d)$ for $T$ timestamps. RE-NET: $O(T \cdot |\mathcal{G}_t|
\cdot d^2)$ where $|\mathcal{G}_t|$ is the size of the temporal subgraph at
time $t$ --- moderate. TiRGN: combines GNN and RNN, $2\times$ cost of RE-NET.
TANGO: ODE solver adds significant overhead ($\approx10\times$ RE-NET) but is
more expressive for continuous dynamics. For KG-Commit: RE-NET is the
practical choice; TANGO is research-grade.
\end{effbox}

\begin{implbox}
\textbf{Prior Implementations and Difficulty.}
\textbf{PyKEEN}: TTransE, TNTComplEx, ATiSE all available.
\textbf{RE-NET}: PyTorch reference code on GitHub (Jin et al.), maintained.
\textbf{TiRGN}: PyTorch reference code (Li et al.).
\textbf{TANGO}: research code.
KG-Commit already has TIME nodes and INTERVAL nodes from the construction
phase --- these map directly to temporal KGE timestamps. Integration requires
attaching timestamps to commit triples in the exported file.
Difficulty: TTransE/TNTComplEx \textbf{easy}; RE-NET \textbf{moderate};
TiRGN \textbf{moderate}; TANGO \textbf{hard}.
\end{implbox}

\begin{prosbox}
\textbf{Pros.}
\textbf{The most natural fit for KG-Commit's commit-sequence structure.}
Bug patterns have strong temporal locality (bug-introducing commits cluster
in time, often preceded by rapid change bursts). RE-NET's subgraph RNN
directly models ``how does the commit neighbourhood evolve over time?''
KG-Commit's TIME and INTERVAL nodes provide structured temporal anchors
for all these methods without extra engineering.
\end{prosbox}

\begin{consbox}
\textbf{Cons.}
Temporal KGE methods are transductive in time: cannot extrapolate beyond
the training time window without retraining or online updates. RE-NET requires
a fixed time granularity (daily? per-commit?), and misalignment with
KG-Commit's variable commit rate must be handled. TANGO's ODE solver is
computationally prohibitive at scale. New or small repositories have
cold-start problems.
\end{consbox}

\paragraph{Key papers.} TTransE (Leblay \& Chekol, WWW 2018), TNTComplEx
(Lacroix et al., ICLR 2020), RE-NET (Jin et al., ICLR 2020), TiRGN (Li et al.,
AAAI 2022), TANGO (Han et al., EMNLP 2021).

%%=============================================================================
\section{Family 11: Subgraph and Inductive Reasoning}
%%=============================================================================

\subsection{What It Captures}

Inductive KG methods reason over local subgraphs around query entities,
enabling generalisation to unseen entities and entirely new graphs (new
repositories). GraIL (Teru et al., 2020) extracts the local enclosing
subgraph around a query triple $(h, r, ?)$ and applies a GNN to classify
whether the triple holds. SNRI (Xu et al., 2022) additionally incorporates
node and relation features. NodePiece (Galkin et al., 2022) represents unseen
entities by their anchor distances, enabling fully inductive KGE.
INDIGO (Liu et al., 2021) uses a differentiable inductive logic system.

\paragraph{Subgraph formulation for KG-Commit.}
\begin{lstlisting}
Query: will COMMIT c touch a FUNCTION f with an existing bug pattern?
Enclosing subgraph: all nodes within k=2 hops of (COMMIT_c, FUNCTION_f)
  Includes: AUTHOR, FILE, co-changed FILEs, called FUNCTIONs, ISSUES
GraIL GNN: classify (COMMIT_c, INTRODUCES_BUG, ?) on enclosing subgraph
Inductive benefit: new repos at test time not seen during training
\end{lstlisting}

\begin{effbox}
\textbf{Efficiency.}
GraIL: enclosing subgraph extraction is $O(k^2 \cdot \bar{d})$ for average
degree $\bar{d}$ and $k$ hops. GNN over the subgraph: $O(|V_\text{sub}|
\cdot d^2)$. For $k=2$ and KG-Commit's average degree $\approx 10$: subgraphs
have $\approx$100--500 nodes --- fast. Training: similar to GNN baselines.
NodePiece: $O(N \cdot K \cdot d)$ for $K$ anchor tokens per entity --- highly
efficient for large KGs.
\end{effbox}

\begin{implbox}
\textbf{Prior Implementations and Difficulty.}
\textbf{GraIL}: PyTorch reference code (Teru et al., ICML 2020), maintained.
\textbf{SNRI}: PyTorch reference code.
\textbf{NodePiece}: PyTorch Geometric integration, maintained.
\textbf{INDIGO}: research-grade.
GraIL requires a subgraph extraction function (implementable in NetworkX or
Cypher in $<$1 day). The inductive setting is critical for KG-Commit: the
model must generalise to new Apache projects not seen during training.
Difficulty: GraIL \textbf{moderate}; NodePiece \textbf{easy}; SNRI
\textbf{moderate}.
\end{implbox}

\begin{prosbox}
\textbf{Pros.}
\textbf{Critical for cross-project generalisation}: GraIL trained on 5 Apache
projects should predict defects in a 6th unseen project without retraining.
NodePiece enables deploying a trained model on any new repository. Enclosing
subgraphs are interpretable --- the extracted subgraph IS the explanation.
SNRI's feature incorporation uses the rich node features from KG-Commit's
construction phase directly.
\end{prosbox}

\begin{consbox}
\textbf{Cons.}
Subgraph extraction adds a preprocessing step per query at inference time.
For dense KGs (high-degree nodes), enclosing subgraphs can be large and
slow to extract. GraIL's original formulation focuses on link prediction,
not node classification --- adapting to commit-level prediction requires
a task reformulation. INDIGO is too immature for production use.
\end{consbox}

\paragraph{Key papers.} GraIL (Teru et al., ICML 2020), SNRI (Xu et al.,
EMNLP 2022), NodePiece (Galkin et al., ICLR 2022), INDIGO (Liu et al.,
NeurIPS 2021).

%%=============================================================================
\section{Family 12: LLM-Augmented Knowledge Graph Inference}
%%=============================================================================

\subsection{What It Captures}

LLM-augmented KG methods combine the symbolic structure of a KG with the
language understanding and world knowledge of large language models.
KnowPrompt (Chen et al., 2022) injects KG entity embeddings into prompt
templates for relation classification. Think-on-Graph (Sun et al., 2024)
has the LLM interactively walk the KG, selecting which relations to follow.
KG-GPT (Kim et al., 2024) prompts GPT-4 with serialised KG subgraphs.
KEPLER (Wang et al., 2021) jointly pretrains a BERT-like LM with KGE
objectives, producing embeddings that are simultaneously language- and
graph-grounded.

\paragraph{Formulation for KG-Commit.}
\begin{lstlisting}
KGLLM prompting approach:
  1. Serialise 2-hop neighbourhood of COMMIT_c as text:
     "Commit c modified F1, F2. F1 has smell GodClass. F2 was
      recently fixed for issue #123. Author a has low ownership."
  2. LLM reasons: "Is this commit likely to introduce a bug?"
  3. Answer grounded by KG facts, not hallucinated

Think-on-Graph:
  LLM iteratively selects next KG hop to follow from COMMIT node
  until sufficient evidence for a prediction is accumulated

KEPLER:
  Fine-tune CodeBERT jointly on code MLM + KG triple scoring
  Use resulting embeddings as node features for R-GCN (Family 4)
\end{lstlisting}

\begin{effbox}
\textbf{Efficiency.}
KnowPrompt: fine-tuning cost of a BERT-class model ($\approx$110M params);
inference $\approx$100ms per query on GPU. KGLLM/KG-GPT: API inference cost
($\sim$\$0.01--\$0.10 per query with GPT-4, depending on context size).
Think-on-Graph: multiple LLM calls per query (iterative KG walking),
$3$--$5\times$ cost of single-pass KGLLM. KEPLER pretraining: significant
GPU cost ($\approx$1--4 GPU-days), but inference is standard BERT speed.
API-based methods are cost-limited, not compute-limited.
\end{effbox}

\begin{implbox}
\textbf{Prior Implementations and Difficulty.}
\textbf{KEPLER}: Hugging Face model checkpoint available.
\textbf{KnowPrompt}: PyTorch reference code (Chen et al., ACL 2022).
\textbf{Think-on-Graph}: PyTorch + LLM API (Sun et al., ICLR 2024).
KGLLM-style prompting is the fastest to prototype (1 day): write a Cypher
query to serialise a commit's neighbourhood as text, pass to Claude/GPT-4.
KEPLER fine-tuning on Apache code commits is a medium-term investment.
Difficulty: KGLLM prompt \textbf{easy}; KnowPrompt \textbf{moderate};
KEPLER fine-tuning \textbf{moderate}; Think-on-Graph \textbf{moderate}.
\end{implbox}

\begin{prosbox}
\textbf{Pros.}
LLMs bring \emph{world knowledge} and \emph{natural language reasoning}
not available to structural GNNs: they understand what ``NPE in concurrent
context'' implies. KGLLM-style grounding reduces hallucination by anchoring
responses to KG facts. Think-on-Graph is highly interpretable --- the
traversal trace IS the explanation. KEPLER produces code + KG-grounded
embeddings that can replace CodeBERT in the GNN node initialisation
pipeline.
\end{prosbox}

\begin{consbox}
\textbf{Cons.}
API-based LLMs have per-query cost prohibitive for online commit analysis
at scale ($10^3$--$10^4$ commits/day). LLM reasoning over serialised KG text
loses graph topology information. KEPLER pretraining is expensive and may
need domain-specific retraining. LLM outputs are non-deterministic and hard
to calibrate. Context window limits restrict the amount of KG neighbourhood
that can be serialised per query.
\end{consbox}

\paragraph{Key papers.} KEPLER (Wang et al., TACL 2021), KnowPrompt (Chen et al.,
ACL 2022), Think-on-Graph (Sun et al., ICLR 2024), KG-GPT (Kim et al.,
EACL 2024).

%%=============================================================================
\section{Family 13: Contrastive and Self-Supervised Representations}
%%=============================================================================

\subsection{What It Captures}

Self-supervised and contrastive KGE methods reduce dependence on scarce defect
labels by learning representations from graph structure alone. SimKGC (Wang et al.,
2022) applies contrastive learning with in-batch negative sampling over
BERT-encoded entity descriptions. KGE-CL (Yu et al., 2021) applies SimCLR-style
augmentations to KG subgraphs (drop edges, mask entities) and trains a GNN with
contrastive loss. GraphMAE (Hou et al., 2022) applies masked autoencoding
(BERT-style) to graph nodes, masking subsets and reconstructing features.
BGRL (Thakoor et al., 2022) is a graph-specific bootstrapped representation
learning method requiring no negative samples.

\paragraph{Formulation for KG-Commit.}
\begin{lstlisting}
Self-supervised pre-training (no defect labels needed):
  Augmentation 1: drop MODIFIES edges with p=0.2
  Augmentation 2: mask FUNCTION node features
  Contrastive loss: maximise agreement between augmented views

GraphMAE: reconstruct masked FUNCTION features from neighbourhood context
BGRL: bootstrap from online/target network pair -- no negatives needed

Transfer: pre-train on full unlabelled KG, then fine-tune on
          the small set of defect-labelled commits
\end{lstlisting}

\begin{effbox}
\textbf{Efficiency.}
SimKGC: BERT-based, $\approx$100ms per batch inference. KGE-CL / BGRL:
GNN-based, $O(|\mathcal{E}| \cdot d^2)$ per epoch with two forward passes
(two augmented views). GraphMAE: adds a decoder network ($O(d^2)$ extra
parameters). BGRL: single forward pass --- the fastest self-supervised GNN;
no negatives required. Pre-training on the unlabelled KG: 1--10 hours;
fine-tuning on defect labels: minutes.
\end{effbox}

\begin{implbox}
\textbf{Prior Implementations and Difficulty.}
\textbf{PyTorch Geometric}: GraphMAE, BGRL, and DGI (Deep Graph Infomax)
all available. \textbf{SimKGC}: reference code (PyTorch, Wang et al., ACL 2022).
For KG-Commit: BGRL is the recommended entry point (simplest, no hyperparameter
sensitivity from negative sampling). GraphMAE requires a feature reconstruction
decoder. SimKGC requires textual entity descriptions (available from commit
messages and Javadoc). Difficulty: BGRL \textbf{easy}; GraphMAE \textbf{easy};
SimKGC \textbf{moderate}; KGE-CL \textbf{moderate}.
\end{implbox}

\begin{prosbox}
\textbf{Pros.}
\textbf{Critical given KG-Commit's label scarcity}: most commits are not
labelled as buggy or clean. Self-supervised pre-training extracts rich
representations from the unlabelled majority. BGRL's bootstrapped approach
avoids the instability of contrastive learning with hard negatives. GraphMAE's
masked-feature pretext task is directly relevant --- reconstructing masked code
metrics from graph context is a strong pretraining signal for defect prediction.
\end{prosbox}

\begin{consbox}
\textbf{Cons.}
Augmentation choices (which edges/features to drop) critically affect quality
and must be tuned for KG-Commit's schema. BGRL may underfit on very
heterogeneous graphs (designed primarily for homogeneous graphs). GraphMAE's
reconstruction quality depends on informative node features. Self-supervised
pre-training adds a two-stage pipeline (pre-train then fine-tune), increasing
total development time. Pre-trained representations may capture structure
unrelated to defect prediction.
\end{consbox}

\paragraph{Key papers.} SimKGC (Wang et al., ACL 2022), GraphMAE (Hou et al.,
KDD 2022), BGRL (Thakoor et al., ICLR 2022), DGI (Veli\v{c}kovi\'{c} et al.,
ICLR 2019).

%%=============================================================================
\section{Family 14: Graph Transformer Architectures}
%%=============================================================================

\subsection{What It Captures}

Graph Transformers apply the self-attention mechanism directly over graph nodes,
using structural information as positional/bias terms. Graphormer (Ying et al.,
2021) encodes spatial distances between nodes as attention biases and uses degree
centrality as node feature augmentation. GPS (Ramp\'{a}\v{s}ek et al., 2022)
combines local MPNN layers with global transformer layers, achieving SOTA on
many graph classification benchmarks. The Structure-Aware Transformer (SAT, Chen
et al., 2022) extracts a subgraph representation for each node and uses it as
the key/value in attention. TokenGT (Kim et al., 2022) treats all nodes and
edges as tokens in a sequence, enabling full self-attention without graph-specific
inductive bias.

\paragraph{Formulation for KG-Commit.}
\begin{lstlisting}
GPS over KG-Commit subgraph:
  Local:  R-GCN message passing over typed edges (Family 4)
  Global: transformer self-attention over all nodes in subgraph
  PE:     Laplacian eigenvectors + random walk structural encoding
  Output: COMMIT node representation -> defect head (sigmoid)

Graphormer:
  Spatial encoding = shortest path distance between all node pairs
  Degree encoding  = learnable bias from in/out degree

TokenGT:
  Flatten all (node, edge) tokens into a sequence, attend globally
\end{lstlisting}

\begin{effbox}
\textbf{Efficiency.}
Graphormer/SAT/TokenGT: $O(N^2 \cdot d)$ per layer due to self-attention
(quadratic in node count). Practical for subgraphs up to $\sim$5000 nodes.
For commit-level subgraphs ($<$1000 nodes), this is tractable on GPU.
GPS: $O(|\mathcal{E}| \cdot d^2 + N^2 \cdot d)$ per layer (MPNN + attention);
the attention term dominates for dense subgraphs. Flash Attention reduces memory
but not FLOPs. Laplacian positional encodings add a one-time preprocessing cost
per graph.
\end{effbox}

\begin{implbox}
\textbf{Prior Implementations and Difficulty.}
\textbf{PyTorch Geometric}: \texttt{GraphormerConv}, GPS layer (via
\texttt{GPSConv}), positional encoding transforms.
\textbf{Graphormer}: reference code (Microsoft Research, production-grade).
\textbf{GPS}: reference code (Ramp\'{a}\v{s}ek et al., NeurIPS 2022), maintained.
\textbf{TokenGT}: reference code (Kim et al., NeurIPS 2022). GPS is the
recommended entry point (best benchmark performance, maintained code, modular
design combining MPNN + transformer). Difficulty: Graphormer \textbf{moderate};
GPS \textbf{moderate}; SAT/TokenGT \textbf{hard} (compute budget issues at scale).
\end{implbox}

\begin{prosbox}
\textbf{Pros.}
Graph Transformers are the current SOTA on many graph classification
benchmarks. GPS's hybrid architecture directly combines the local relational
reasoning of R-GCN (Family~\ref{sec:rgcn}) with the global context of
transformers --- ideal for KG-Commit where both local code structure and
global project-level patterns matter. Graphormer's spatial encoding captures
graph-distance information explicitly, relevant for CPG node distances.
\end{prosbox}

\begin{consbox}
\textbf{Cons.}
$O(N^2)$ self-attention is expensive for large graphs. Positional encodings
(Laplacian eigenvectors) are sensitive to graph changes and add preprocessing
cost. Flash Attention reduces memory but not FLOPs. The performance gains over
well-tuned GNNs are modest on many tasks (GPS beats GCN by 3--5\% on most
benchmarks). Not widely tested on code/KG tasks. SAT and TokenGT are research-
grade for large heterogeneous graphs.
\end{consbox}

\paragraph{Key papers.} Graphormer (Ying et al., NeurIPS 2021), SAT (Chen et al.,
NeurIPS 2022), GPS (Ramp\'{a}\v{s}ek et al., NeurIPS 2022), TokenGT (Kim et al.,
NeurIPS 2022).

%%=============================================================================
\section{Comprehensive Comparison}
%%=============================================================================

\begin{longtable}{p{3.0cm}p{1.3cm}p{1.3cm}p{1.2cm}p{1.2cm}p{1.7cm}p{1.8cm}}
\toprule
\textbf{Family} & \textbf{Hetero.} & \textbf{Temp.} & \textbf{Induct.} & \textbf{Interp.} & \textbf{Cost} & \textbf{KG-Commit fit} \\
\midrule
\endhead
Shallow GCN/GraphSAGE  & Partial    & --         & \checkmark & --         & Very low  & Baseline \\
Attention (GAT/HAN)    & \checkmark & --         & Partial    & \checkmark & Low       & Good \\
Gated GNN (GGNN)       & Partial    & --         & --         & --         & Low       & CPG-specific \\
Relational (R-GCN/HGT) & \checkmark & --         & Partial    & --         & Medium    & \textbf{Best structural} \\
TransE family          & --         & --         & --         & --         & Very low  & Link prediction \\
Bilinear (DistMult/ComplEx) & --    & --         & --         & --         & Very low  & Link prediction \\
Rotation (RotatE/HAKE) & --         & --         & --         & Partial    & Low       & \textbf{Best KGE} \\
Path reasoning (PRA/MINERVA) & \checkmark & --   & Partial    & \checkmark & Medium    & Explainability \\
Rule learning (AMIE/Neural-LP) & \checkmark & -- & \checkmark & \checkmark & Low       & \textbf{Interpretability} \\
Temporal (TTransE/RE-NET) & Partial & \checkmark & --         & --         & Medium    & \textbf{Temporal signal} \\
Subgraph (GraIL/NodePiece) & \checkmark & --    & \checkmark & \checkmark & Medium    & \textbf{Cross-project} \\
LLM-augmented (KGLLM)  & \checkmark & --         & \checkmark & \checkmark & High      & NL grounding \\
Contrastive (BGRL/GraphMAE) & Partial & --       & \checkmark & --         & Medium    & \textbf{Label scarcity} \\
Graph Transformer (GPS) & \checkmark & --        & Partial    & --         & High      & Best accuracy \\
\bottomrule
\end{longtable}

%%=============================================================================
\section{Recommended Phased Inference Plan}
%%=============================================================================

\begin{enumerate}[leftmargin=1.5em]
  \item \textbf{Phase 1 --- KGE Baseline (Week 1--2).}
        Export KG-Commit triples and train TransE and RotatE via PyKEEN.
        Use for link prediction (COMMIT$\to$INTRODUCES\_BUG). Establish a
        strong embedding baseline with zero GNN infrastructure.

  \item \textbf{Phase 2 --- Shallow GNN Baseline.}
        Build a PyG \texttt{HeteroData} object from the KG. Train GraphSAGE
        (inductive) with code metric node features. This is the JIT-DP
        baseline against which all structural improvements are measured.

  \item \textbf{Phase 3 --- R-GCN / HGT (Week 4--6).}
        Add R-GCN with basis decomposition over all edge types. Tune number
        of layers (2--3), hidden dimension (128--256). Compare against the
        shallow GNN. This is the primary structural model.

  \item \textbf{Phase 4 --- GGNN on CPG Subgraphs.}
        For commits with CPG subgraphs (Joern output), apply GGNN as a local
        feature extractor. Feed its output as a COMMIT node attribute into
        the R-GCN.

  \item \textbf{Phase 5 --- Self-Supervised Pre-Training.}
        Pre-train R-GCN or HGT with BGRL or GraphMAE on the unlabelled
        portion of the KG. Fine-tune on defect labels. Evaluate benefit for
        label-scarce scenarios.

  \item \textbf{Phase 6 --- Temporal Extension.}
        Add TNTComplEx or RE-NET to incorporate commit sequence information.
        Evaluate whether temporal signal adds value over structural signal
        alone.

  \item \textbf{Phase 7 --- Rule Mining.}
        Run AMIE on the KG triple file. Review mined rules with domain experts.
        Encode actionable rules as code review guidelines. Integrate
        high-confidence rules as hard constraints in the GNN.

  \item \textbf{Phase 8 --- Cross-Project / Inductive.}
        Apply GraIL or NodePiece for inductive evaluation on held-out Apache
        projects. This evaluates the model's generalisation beyond the training
        corpus.

  \item \textbf{Phase 9 --- LLM Grounding (Optional).}
        Prototype KGLLM-style prompting for a small set of high-confidence
        defect predictions. Use as an explainability layer on top of the GNN
        output, not as the primary predictor.

  \item \textbf{Phase 10 --- Graph Transformer.}
        Replace R-GCN with GPS as the backbone. Compare accuracy vs.\ training
        cost. Use only if the accuracy gain justifies the compute overhead.
\end{enumerate}

%%=============================================================================
\section{Open Research Questions}
%%=============================================================================

\begin{itemize}[leftmargin=1.5em]
  \item \textbf{Heterogeneous GNN for JIT-DP.} No published study directly
        applies HGT or R-GCN to a multi-typed code KG for defect prediction.
        This is a core open gap that KG-Commit addresses.
  \item \textbf{Temporal GNN for commit streams.} RE-NET and TiRGN were
        evaluated on event-based KGs (ICEWS, GDELT) --- their transferability
        to commit-level code KGs is untested.
  \item \textbf{Label-efficient learning.} Can self-supervised pre-training
        (BGRL/GraphMAE) on the unlabelled KG reduce the labelled-data
        requirement below the current literature minimum ($\sim$1000 defective
        commits)?
  \item \textbf{Rule mining from code KGs.} AMIE has not been applied to
        commit-level code KGs. What defect-predictive rules can be mined, and
        are they project-portable?
  \item \textbf{Subgraph reasoning for cross-project transfer.} GraIL's
        inductive framework enables zero-shot evaluation on new projects. What
        is the accuracy gap between in-distribution and cross-project settings
        for KG-Commit?
  \item \textbf{Composing GNN and KGE.} How to best combine structural GNN
        representations (R-GCN) with relational KGE representations (RotatE)
        in a unified model for code-change defect prediction?
  \item \textbf{Explainability for commit review.} Which inference method
        produces the most actionable explanation for a developer reviewing a
        flagged commit --- paths, rules, attention weights, or LLM traces?
\end{itemize}

%%=============================================================================
\section*{References}
%%=============================================================================

\begin{enumerate}[leftmargin=1.5em, label={[\arabic*]}]
  \item T. Kipf and M. Welling, ``Semi-Supervised Classification with Graph
        Convolutional Networks,'' \textit{ICLR}, 2017.
  \item W. Hamilton et al., ``Inductive Representation Learning on Large
        Graphs,'' \textit{NeurIPS}, 2017.
  \item M. Defferrard et al., ``Convolutional Neural Networks on Graphs with
        Fast Localized Spectral Filtering,'' \textit{NeurIPS}, 2016.
  \item F. Wu et al., ``Simplifying Graph Convolutional Networks,''
        \textit{ICML}, 2019.
  \item P. Veli\v{c}kovi\'{c} et al., ``Graph Attention Networks,''
        \textit{ICLR}, 2018.
  \item S. Brody et al., ``How Attentive are Graph Attention Networks?''
        \textit{ICLR}, 2022.
  \item X. Wang et al., ``Heterogeneous Graph Attention Network,''
        \textit{WWW}, 2019.
  \item Y. Li et al., ``Gated Graph Sequence Neural Networks,'' \textit{ICLR},
        2016.
  \item Z. Zhou et al., ``Devign: Effective Vulnerability Identification by
        Learning Comprehensive Program Semantics via Graph Neural Networks,''
        \textit{NeurIPS}, 2019.
  \item K. Tai et al., ``Improved Semantic Representations from Tree-Structured
        Long Short-Term Memory Networks,'' \textit{ACL}, 2015.
  \item M. Schlichtkrull et al., ``Modeling Relational Data with Graph
        Convolutional Networks,'' \textit{ESWC}, 2018.
  \item Z. Hu et al., ``Heterogeneous Graph Transformer,'' \textit{WWW}, 2020.
  \item S. Vashishth et al., ``Composition-based Multi-Relational Graph
        Convolutional Networks,'' \textit{ICLR}, 2020.
  \item A. Bordes et al., ``Translating Embeddings for Modeling Multi-relational
        Data,'' \textit{NeurIPS}, 2013.
  \item Z. Wang et al., ``Knowledge Graph Embedding by Translating on
        Hyperplanes,'' \textit{AAAI}, 2014.
  \item Y. Lin et al., ``Learning Entity and Relation Embeddings for Knowledge
        Graph Completion,'' \textit{AAAI}, 2015.
  \item B. Yang et al., ``Embedding Entities and Relations for Learning and
        Inference in Knowledge Bases,'' \textit{ICLR}, 2015.
  \item T. Trouillon et al., ``Complex Embeddings for Simple Link Prediction,''
        \textit{ICML}, 2016.
  \item I. Bala\v{z}evi\'{c} et al., ``TuckER: Tensor Factorization for
        Knowledge Graph Completion,'' \textit{EMNLP}, 2019.
  \item M. Nickel et al., ``A Three-Way Model for Collective Learning on
        Multi-Relational Data,'' \textit{ICML}, 2011.
  \item Z. Sun et al., ``RotatE: Knowledge Graph Embedding by Relational
        Rotation in Complex Space,'' \textit{ICLR}, 2019.
  \item Z. Zhang et al., ``Learning Hierarchy-Aware Knowledge Graph Embeddings
        for Link Prediction,'' \textit{ACL}, 2020.
  \item L. Chao et al., ``PairRE: Knowledge Graph Embeddings via Paired
        Relation Vectors,'' \textit{ACL}, 2021.
  \item S. Zhang et al., ``Quaternion Knowledge Graph Embeddings,''
        \textit{NeurIPS}, 2019.
  \item N. Lao et al., ``Random Walk Inference and Learning in a Large Scale
        Knowledge Base,'' \textit{EMNLP}, 2011.
  \item W. Xiong et al., ``DeepPath: A Reinforcement Learning Method for
        Knowledge Graph Reasoning,'' \textit{EMNLP}, 2017.
  \item R. Das et al., ``Go for a Walk and Arrive at the Answer: Reasoning
        Over Paths in Knowledge Bases,'' \textit{ICLR}, 2018.
  \item Y. Shen et al., ``M-Walk: Learning to Walk over Graphs using Monte
        Carlo Tree Search,'' \textit{NeurIPS}, 2018.
  \item L. Gal\'{a}rraga et al., ``AMIE: Association Rule Mining under
        Incomplete Evidence in Ontological Knowledge Bases,'' \textit{WWW},
        2013.
  \item F. Yang et al., ``Differentiable Learning of Logical Rules for
        Knowledge Base Reasoning,'' \textit{NeurIPS}, 2017.
  \item A. Sadeghian et al., ``DRUM: End-To-End Differentiable Rule Mining
        on Knowledge Graphs,'' \textit{NeurIPS}, 2019.
  \item M. Qu et al., ``RNNLogic: Learning Logic Rules for Reasoning on
        Knowledge Graphs,'' \textit{ICLR}, 2021.
  \item J. Leblay and M. Chekol, ``Deriving Validity Time in Knowledge Graph,''
        \textit{WWW}, 2018.
  \item T. Lacroix et al., ``Tensor Decompositions for Temporal Knowledge
        Base Completion,'' \textit{ICLR}, 2020.
  \item W. Jin et al., ``Recurrent Event Network: Autoregressive Structure
        Inference over Temporal Knowledge Graphs,'' \textit{ICLR}, 2020.
  \item Z. Li et al., ``TiRGN: Time-Guided Recurrent Graph Network with
        Local-Global Historical Patterns for Temporal Knowledge Graph
        Reasoning,'' \textit{AAAI}, 2022.
  \item Z. Han et al., ``TANGO: Time-Reversal Latent GraphODE for
        Multi-Relational Temporal Knowledge Graph Completion,''
        \textit{EMNLP}, 2021.
  \item K. Teru et al., ``Inductive Relation Prediction by Subgraph
        Reasoning,'' \textit{ICML}, 2020.
  \item M. Xu et al., ``Subgraph Neighboring Relations Infomax for Inductive
        Link Prediction on Knowledge Graphs,'' \textit{EMNLP}, 2022.
  \item M. Galkin et al., ``NodePiece: Compositional and Parameter-Efficient
        Representations of Large Knowledge Graphs,'' \textit{ICLR}, 2022.
  \item Z. Liu et al., ``INDIGO: GNN-Based Inductive Knowledge Graph
        Completion Using Pair-Wise Encoding,'' \textit{NeurIPS}, 2021.
  \item X. Wang et al., ``KEPLER: A Unified Model for Knowledge Embedding
        and Pre-trained Language Representation,'' \textit{TACL}, 2021.
  \item Z. Chen et al., ``KnowPrompt: Knowledge-aware Prompt-tuning with
        Synergistic Optimization for Relation Extraction,'' \textit{ACL}, 2022.
  \item J. Sun et al., ``Think-on-Graph: Deep and Responsible Reasoning of
        LLMs with Knowledge Graph,'' \textit{ICLR}, 2024.
  \item W. Kim et al., ``KG-GPT: A General Framework for Reasoning on
        Knowledge Graphs Using Large Language Models,'' \textit{EACL}, 2024.
  \item W. Wang et al., ``SimKGC: Simple Contrastive Knowledge Graph
        Completion with Pre-trained Language Models,'' \textit{ACL}, 2022.
  \item Z. Hou et al., ``GraphMAE: Self-Supervised Masked Graph Autoencoders,''
        \textit{KDD}, 2022.
  \item S. Thakoor et al., ``Large-Scale Representation Learning on Graphs
        via Bootstrapping,'' \textit{ICLR}, 2022.
  \item P. Veli\v{c}kovi\'{c} et al., ``Deep Graph Infomax,'' \textit{ICLR},
        2019.
  \item C. Ying et al., ``Do Transformers Really Perform Bad for Graph
        Representation?'' \textit{NeurIPS}, 2021.
  \item D. Chen et al., ``Structure-Aware Transformer for Graph
        Representation Learning,'' \textit{NeurIPS}, 2022.
  \item L. Ramp\'{a}\v{s}ek et al., ``Recipe for a General, Powerful, Scalable
        Graph Transformer,'' \textit{NeurIPS}, 2022.
  \item J. Kim et al., ``Pure Transformers are Powerful Graph Learners
        (TokenGT),'' \textit{NeurIPS}, 2022.
  \item M. Ali et al., ``PyKEEN 1.0: A Python Library for Training and
        Evaluating Knowledge Graph Embeddings,'' \textit{JMLR}, 2021.
\end{enumerate}

\end{document}
